% ============================================================================
% Chuong 1: TONG QUAN
% ============================================================================

\chapter{TỔNG QUAN}
\label{ch:overview}

\section{Bối cảnh kiến trúc: từ một khối đến vi dịch vụ trên Kubernetes}
\label{sec:arch-context}

Phần này làm rõ \emph{vì sao} bài toán
giám sát an ninh trong đồ án lại đặt ra trên môi trường Kubernetes chứ không phải trên
hạ tầng truyền thống, và \emph{vì sao} cách quan sát phải đặt ở tầng nhân. Lý do nằm ở
sự dịch chuyển kiến trúc phần mềm trong một thập kỷ qua: từ kiến trúc \emph{một khối}
(\textit{monolithic}) chạy trên máy chủ vật lý, sang kiến trúc \emph{vi dịch vụ}
(\textit{microservices}) đóng gói trong container và điều phối bởi Kubernetes. Sự dịch
chuyển này không chỉ thay đổi cách phần mềm được xây dựng và mở rộng mà còn làm thay đổi
căn bản \emph{bề mặt tấn công}, khiến các biện pháp an ninh truyền thống trở nên kém hiệu
quả --- đó chính là khoảng trống mà đồ án hướng tới.

\subsection{Kiến trúc một khối và mở rộng theo chiều dọc}

Trong kiến trúc một khối, toàn bộ chức năng của ứng dụng --- từ tầng giao diện, tầng
nghiệp vụ đến tầng truy cập dữ liệu --- được biên dịch và triển khai như \emph{một tiến
trình duy nhất} chạy trên một máy chủ (vật lý hoặc máy ảo) như minh họa ở
Hình~\ref{fig:arch-monolith}. Đặc điểm quan trọng nhất xét từ góc độ an ninh là: các lời
gọi giữa các tầng diễn ra \emph{trong cùng một tiến trình} (in-process), tức không đi qua
mạng. Hệ quả là gần như toàn bộ luồng dữ liệu nội bộ của ứng dụng nằm gọn bên trong một
máy chủ và không hề xuất hiện trên dây mạng để có thể bị quan sát. Khi cần phục vụ nhiều
người dùng hơn, mô hình này mở rộng \emph{theo chiều dọc} (\textit{vertical scaling}):
nâng cấp phần cứng (thêm CPU, RAM) cho chính máy chủ đó. Cách mở rộng này đơn giản nhưng
có trần vật lý và tạo ra một điểm hỏng tập trung.

Trong mô hình này, an ninh được tổ chức theo nguyên lý \emph{phòng thủ chu vi}
(\textit{perimeter security}): một tường lửa hoặc hệ thống phát hiện xâm nhập đặt ở
\emph{biên mạng}, kiểm soát lưu lượng ra vào giữa Internet và máy chủ --- thường gọi là
lưu lượng \emph{Bắc--Nam}. Vì luồng nội bộ là in-process, mô hình phòng thủ chu vi này
là \emph{đủ}: kẻ tấn công muốn tới được dữ liệu phải vượt qua biên, và biên là nơi duy
nhất cần canh gác.

\begin{figure}[H]
  \centering
  \resizebox{0.82\linewidth}{!}{% Kien truc mot khoi (monolithic) truyen thong + mo rong theo chieu doc. Chuong 2.
\begin{tikzpicture}[font=\footnotesize, >={Stealth[length=2.2mm]},
    node distance=0.7cm and 1.1cm,
    cli/.style ={draw, rounded corners=2pt, fill=blue!8, align=center, minimum height=0.7cm},
    fw/.style  ={draw, fill=red!12, align=center, minimum width=1.5cm, minimum height=1.6cm},
    mod/.style ={draw, rounded corners=2pt, fill=orange!16, align=center,
                 minimum width=4.2cm, minimum height=0.72cm},
    db/.style  ={draw, cylinder, shape border rotate=90, aspect=0.28, fill=green!14,
                 minimum width=1.7cm, minimum height=1.0cm, align=center},
  ]
  \node[cli] (cli) {Người dùng\\(Internet)};
  \node[fw, right=of cli] (fw) {Tường lửa\\chu vi\\(Bắc--Nam)};
  % Cac tang cua monolith
  \node[mod, right=2.0cm of fw, yshift=1.45cm] (m1) {Tầng giao diện};
  \node[mod, below=0.32cm of m1] (m2) {Tầng nghiệp vụ};
  \node[mod, below=0.32cm of m2] (m3) {Tầng truy cập dữ liệu};
  \node[db, below=0.5cm of m3] (db) {CSDL};
  \draw[<->, gray!70] (m1) -- (m2);
  \draw[<->, gray!70] (m2) -- (m3);
  \draw[<->, gray!70] (m3) -- (db);
  % Khung host (KHONG to nen de khong de len noi dung)
  \node[draw, thick, rounded corners=3pt, fit=(m1)(m2)(m3)(db), inner sep=10pt,
        label={[font=\bfseries]above:Máy chủ (VM / vật lý) --- một tiến trình}] (host) {};
  \draw[->] (cli) -- (fw);
  \draw[->] (fw.east) -- (host.west|-fw.east);
  % Chu thich
  \node[align=left, font=\scriptsize, below=0.2cm of host.south, anchor=north, text width=7.2cm]
    {\textbf{Lời gọi giữa các tầng = trong tiến trình} (in-process), \emph{không} đi qua mạng
     $\Rightarrow$ không có lưu lượng nội bộ để giám sát ở tầng mạng.};
  % Mo rong doc
  \draw[->, blue!60, line width=1pt] ($(host.north east)+(0.4,0)$) -- ++(0,-2.7)
    node[midway, right, align=left, font=\scriptsize] {Mở rộng\\\textbf{theo chiều dọc}:\\thêm CPU/RAM\\cho cùng máy};
\end{tikzpicture}
}
  \caption{Kiến trúc một khối (monolith) truyền thống.}
  \label{fig:arch-monolith}
\end{figure}

\subsection{Kiến trúc vi dịch vụ trên Kubernetes và mở rộng theo chiều ngang}

Kiến trúc vi dịch vụ phân rã ứng dụng thành nhiều dịch vụ nhỏ, độc lập, mỗi dịch vụ đóng
gói trong một container và được điều phối bởi Kubernetes \cite{berardi2022microservice}.
Như Hình~\ref{fig:arch-micro} thể hiện, các dịch vụ này chạy phân tán trên nhiều nút
(\textit{node}) của cụm và giao tiếp với nhau \emph{qua mạng} thông qua một mạng phủ
(\textit{overlay network}) phẳng. Mô hình mở rộng \emph{theo chiều ngang} (\textit{horizontal
scaling}): thay vì nâng cấp một máy, hệ thống nhân bản thêm các bản sao pod hoặc thêm nút
vào cụm. Điều này mang lại độ co giãn và khả năng chịu lỗi cao hơn hẳn, là lý do kiến trúc
này trở thành tiêu chuẩn thực tế cho các hệ thống quy mô lớn hiện đại.

Tuy nhiên, chính sự phân tán này sinh ra một loại lưu lượng mới mà mô hình một khối không
có: lưu lượng \emph{Đông--Tây} --- các lời gọi dịch vụ-tới-dịch vụ (pod-tới-pod) đi
\emph{qua mạng} bên trong cụm. Trong một hệ thống vi dịch vụ thực tế, lưu lượng Đông--Tây
thường lớn hơn nhiều lần lưu lượng Bắc--Nam. Mỗi pod giờ đây vừa là một điểm có thể bị
khai thác, vừa là một bàn đạp để kẻ tấn công \emph{dịch chuyển ngang} (\textit{lateral
movement}) sang các pod khác. Vòng đời pod ngắn và địa chỉ IP thay đổi liên tục càng làm
cho việc theo dõi trở nên khó khăn.

\begin{figure}[H]
  \centering
  \resizebox{0.92\linewidth}{!}{% Kien truc vi dich vu tren Kubernetes (cloud-native) + mo rong ngang + quan sat eBPF. Chuong 2.
\begin{tikzpicture}[font=\footnotesize, >={Stealth[length=2.2mm]},
    node distance=0.6cm and 0.9cm,
    cli/.style ={draw, rounded corners=2pt, fill=blue!8, align=center, minimum height=0.7cm},
    ing/.style ={draw, fill=red!12, align=center, minimum width=1.6cm, minimum height=1.3cm},
    pod/.style ={draw, rounded corners=2pt, fill=orange!16, align=center,
                 minimum width=1.5cm, minimum height=0.72cm},
    dbp/.style ={draw, rounded corners=2pt, fill=green!14, align=center,
                 minimum width=1.5cm, minimum height=0.72cm},
    col/.style ={draw, rounded corners=2pt, fill=red!10, align=center,
                 minimum width=3.6cm, minimum height=0.55cm, font=\scriptsize},
    zone/.style={draw, dashed, rounded corners=3pt},
  ]
  \node[cli] (cli) {Người dùng\\(Internet)};
  \node[ing, right=of cli] (ing) {Ingress /\\Cân bằng tải\\(Bắc--Nam)};
  % Node 1
  \node[pod, right=2.0cm of ing, yshift=1.7cm] (a) {svc A};
  \node[pod, right=0.7cm of a] (b) {svc B};
  \node[col, below=0.3cm of a, xshift=0.8cm] (c1) {collector eBPF (DaemonSet)};
  \node[zone, fit=(a)(b)(c1), inner sep=7pt, label={[font=\scriptsize\bfseries]above:Node 1}] (n1) {};
  % Node 2
  \node[pod, below=1.7cm of a] (c) {svc C};
  \node[dbp, right=0.7cm of c] (d) {CSDL};
  \node[col, below=0.3cm of c, xshift=0.8cm] (c2) {collector eBPF (DaemonSet)};
  \node[zone, fit=(c)(d)(c2), inner sep=7pt, label={[font=\scriptsize\bfseries]below:Node 2}] (n2) {};
  % Cum K8s (khung ngoai, khong to nen)
  \node[draw, thick, rounded corners=4pt, fit=(n1)(n2), inner sep=12pt,
        label={[font=\bfseries]above:Cụm Kubernetes (mạng overlay phẳng)}] (cluster) {};
  % Bac-Nam
  \draw[->] (cli) -- (ing);
  \draw[->] (ing.east) -- (a.west);
  % Dong-Tay (qua mang overlay)
  \draw[<->, purple!70, line width=0.9pt] (a) -- (b);
  \draw[<->, purple!70, line width=0.9pt] (b.south) to[bend left=15] (d.north);
  \draw[<->, purple!70, line width=0.9pt] (c) -- (d);
  \draw[<->, purple!70, line width=0.9pt] (a.south) to[bend right=12] (c.north);
  \node[purple!70, font=\scriptsize, align=center] at ($(cluster.east)+(1.55,0.55)$)
    {lưu lượng\\\textbf{Đông--Tây}\\(qua mạng)};
  % Mo rong ngang
  \draw[->, blue!60, line width=1pt] ($(cluster.east)+(0.4,-0.4)$) -- ++(0,-2.0)
    node[midway, right, align=left, font=\scriptsize] {Mở rộng\\\textbf{theo chiều ngang}:\\nhân bản pod /\\thêm node};
\end{tikzpicture}
}
  \caption{Kiến trúc vi dịch vụ trên Kubernetes.}
  \label{fig:arch-micro}
\end{figure}

\subsection{Quét cổng và dịch chuyển ngang: định nghĩa và mô hình phát hiện}
\label{subsec:threat-defs}

Trên bề mặt Đông--Tây vừa mô tả, hai hành vi tấn công là trọng tâm của đồ án. Mục này
định nghĩa chúng một cách chính xác và --- quan trọng hơn --- chỉ ra \emph{tín hiệu} mà
mỗi hành vi để lại, vì chính các tín hiệu này quyết định giải thuật phát hiện được dùng ở
các mục sau.

\paragraph{Quét cổng (\textit{port scanning}).} Là hành vi một nguồn lần lượt thử kết nối
tới nhiều cổng (quét \emph{dọc} trên một đích) hoặc nhiều đích (quét \emph{ngang} trên một
cổng) nhằm lập bản đồ các dịch vụ đang mở. Trong cụm Kubernetes, đây thường là bước
\emph{trinh sát} đầu tiên sau khi kẻ tấn công kiểm soát được một pod: từ pod đó, chúng dò
quét các pod và dịch vụ lân cận để tìm điểm khai thác tiếp theo. Dấu hiệu đặc trưng của
quét cổng, nhìn từ tầng nhân, là \emph{một nguồn tạo ra số lượng kết nối tăng vọt bất
thường} trong một khoảng thời gian ngắn (Hình~\ref{fig:threat-patterns}, trái). Tín hiệu
này có bản chất \emph{tần suất}, và đó chính là lý do hệ thống dùng cấu trúc ước lượng tần
suất Count-Min Sketch (Mục~\ref{sec:theory-cms}) để bắt nó với chi phí bộ nhớ thấp.

\paragraph{Dịch chuyển ngang (\textit{lateral movement}).} Là giai đoạn kẻ tấn công, sau
khi chiếm một pod, lần lượt mở rộng quyền kiểm soát sang các pod khác để tiến gần tới mục
tiêu giá trị cao (ví dụ cơ sở dữ liệu). Khác với quét cổng vốn lộ ra ở \emph{tần suất},
dịch chuyển ngang lộ ra ở \emph{cấu trúc} của lưu lượng: xuất hiện những chuỗi hoặc chu
trình liên lạc pod-tới-pod bất thường (Hình~\ref{fig:threat-patterns}, phải). Một nhóm pod
liên lạc vòng tròn với nhau là một dấu hiệu đáng ngờ, và phạm vi mà một pod bị chiếm có
thể với tới trong vài bước cho biết mức độ nghiêm trọng của sự cố.

\begin{figure}[H]
  \centering
  \resizebox{0.92\linewidth}{!}{% Hai mau hanh vi: (trai) quet cong = 1 nguon -> nhieu dich; (phai) dich chuyen ngang = chuoi pod. Chuong 2.
\begin{tikzpicture}[font=\footnotesize, >={Stealth[length=2mm]},
    src/.style ={draw, circle, fill=red!22, draw=red!70!black, minimum size=0.85cm,
                 font=\scriptsize\bfseries},
    tgt/.style ={draw, rounded corners=2pt, fill=blue!10, minimum width=1.05cm,
                 minimum height=0.55cm, font=\scriptsize},
    comp/.style={draw, rounded corners=2pt, fill=red!16, draw=red!70!black,
                 minimum width=1.05cm, minimum height=0.6cm, font=\scriptsize\bfseries},
    scan/.style={->, red!70!black, thick},
    lat/.style ={->, red!70!black, very thick},
  ]
  % ===== Panel trai: QUET CONG =====
  \node[src] (s) at (0,0) {S};
  \foreach \i/\y in {1/2.0, 2/1.0, 3/0.0, 4/-1.0, 5/-2.0}{
    \node[tgt] (t\i) at (2.6,\y) {cổng \i};
    \draw[scan] (s) -- (t\i);
  }
  \node[align=center, font=\scriptsize, below=2.7cm of s, yshift=0.3cm] (capL)
    {\textbf{Quét cổng}\\một nguồn $S$ tạo \emph{nhiều} kết nối\\$\Rightarrow$ tín hiệu: \emph{tần suất} cao (CMS)};

  % ===== Panel phai: DICH CHUYEN NGANG =====
  \begin{scope}[xshift=7.2cm]
    \node[comp] (a) at (0,1.0)   {pod A};
    \node[comp] (b) at (1.9,0.0) {pod B};
    \node[comp] (c) at (3.8,1.0) {pod C};
    \node[comp] (d) at (2.9,-1.4){pod D};
    \draw[lat] (a) -- (b);
    \draw[lat] (b) -- (c);
    \draw[lat] (b) -- (d);
    \draw[lat] (c) to[bend left=20] (a);  % chu trinh -> SCC
    \node[align=center, font=\scriptsize, below=0.4cm of d] (capR)
      {\textbf{Dịch chuyển ngang}\\chuỗi/chu trình pod bị chiếm\\$\Rightarrow$ tín hiệu: \emph{cấu trúc đồ thị} (Tarjan/BFS)};
  \end{scope}

  % vach ngan giua hai panel
  \draw[gray!45, dashed] (4.4,2.4) -- (4.4,-3.0);
\end{tikzpicture}
}
  \caption{Hai mẫu hành vi tấn công và tín hiệu tương ứng.}
  \label{fig:threat-patterns}
\end{figure}

\paragraph{Mô hình phát hiện thống nhất.} Hai tín hiệu trên được nắm bắt trên cùng một
biểu diễn: mô hình hóa lưu lượng nội cụm thành một \emph{đồ thị truyền thông} có hướng
$G=(V,E)$, trong đó mỗi đỉnh là một địa chỉ (pod, dịch vụ, nút) và mỗi cung là một kết nối
đã quan sát được. Trên biểu diễn này, bài toán phát hiện tách thành hai nhánh bổ trợ:
(i) \emph{nhánh tần suất} --- ước lượng số kết nối phát ra từ mỗi đỉnh nguồn để phát hiện
quét cổng, giải bằng Count-Min Sketch; và (ii) \emph{nhánh cấu trúc} --- phân tích quan hệ
giữa các đỉnh để phát hiện dịch chuyển ngang, giải bằng giải thuật Tarjan (tìm thành phần
liên thông mạnh) và duyệt BFS theo $k$ bước (ước lượng bán kính ảnh hưởng), trình bày ở
Mục~\ref{sec:graph-ip}. Việc tách bài toán theo hai loại tín hiệu như vậy là kim chỉ nam cho phần cơ sở giải thuật ở Chương~\ref{ch:foundations}.

\subsection{Hệ quả an ninh và lý do quan sát ở tầng nhân}
\label{subsec:why-kernel}

Đặt hai kiến trúc cạnh nhau, có thể thấy ngay vì sao phòng thủ chu vi truyền thống không
còn đủ. Tường lửa biên mạng chỉ thấy lưu lượng Bắc--Nam và gần như \emph{mù} trước toàn bộ
lưu lượng Đông--Tây nội cụm --- đúng nơi mà hành vi quét cổng và dịch chuyển ngang diễn ra.
Khi kẻ tấn công đã chiếm được một pod, chúng có thể tự do dò quét và lan sang các pod khác
mà không một lần nào phải đi qua biên. Đây là lý do mô hình an ninh hiện đại chuyển sang
nguyên lý \emph{Zero-Trust} \cite{nist800207}: không mặc nhiên tin tưởng lưu lượng nội bộ,
phải quan sát và kiểm soát \emph{mọi} kết nối kể cả bên trong vành đai. Các báo cáo công
nghiệp gần đây xác nhận rủi ro này: cấu hình sai và thiếu khả năng quan sát nội cụm là
nhóm nguyên nhân hàng đầu của sự cố an ninh Kubernetes
\cite{redhat2024k8ssec,wiz2025k8ssec,pereiravale2021security}.

Vấn đề kế tiếp là \emph{quan sát ở đâu}. Có ba lựa chọn quen thuộc, mỗi lựa chọn một nhược
điểm trong môi trường vi dịch vụ: (i) cài tác nhân giám sát vào \emph{từng ứng dụng} ---
xâm lấn, phụ thuộc ngôn ngữ lập trình và phải sửa từng dịch vụ; (ii) chèn một
\textit{sidecar} proxy cạnh mỗi pod --- tốn tài nguyên và chỉ thấy lưu lượng tầng ứng dụng
đã được định tuyến qua nó; (iii) phản chiếu (\textit{mirror}) toàn bộ gói tin trên mạng ---
chi phí lớn và khó truy ngữ cảnh tiến trình. Cả ba đều không lý tưởng. Trong khi đó, có một
điểm quan sát \emph{tự nhiên} mà mọi kết nối đều phải đi qua: \emph{nhân của hệ điều hành}
trên mỗi nút. Vì mọi kết nối TCP của mọi pod đều được xử lý bởi nhân của nút chứa pod đó,
một chương trình đặt tại tầng nhân sẽ quan sát được trọn vẹn lưu lượng Đông--Tây mà không
cần sửa ứng dụng, không phụ thuộc ngôn ngữ và chỉ cần hiện diện một lần trên mỗi nút. Công
nghệ cho phép làm điều đó một cách an toàn và hiệu năng cao chính là \textbf{eBPF}, được trình bày ở Chương~\ref{ch:foundations}. Như vậy, lựa chọn quan sát tầng nhân không phải ngẫu nhiên
mà là hệ quả trực tiếp của đặc thù kiến trúc vi dịch vụ trên Kubernetes.

\section{Khảo sát nhu cầu thực tế tại doanh nghiệp}
\label{sec:survey-enterprise}

Đề tài không xuất phát từ một giả định lý thuyết thuần túy mà bắt nguồn từ một nhu cầu
vận hành có thật, được ghi nhận qua quá trình khảo sát tại một doanh nghiệp đang vận hành
nền tảng Kubernetes nội bộ. Mục này trình bày môi trường khảo sát, phương pháp thu thập
yêu cầu, hiện trạng giám sát an ninh và các điểm đau, từ đó tổng hợp thành tập nhu cầu
làm đầu vào cho phần phân tích yêu cầu (Mục~\ref{sec:requirements}). Việc hiểu rõ bài
toán thực tế là cơ sở để bảo đảm hệ thống được thiết kế đúng nhu cầu thay vì xuất phát
từ giả định.

Về phương pháp luận, đồ án tuân theo quy trình \emph{kỹ nghệ yêu cầu}
(\textit{requirements engineering}) gồm bốn bước nối tiếp: (i) \emph{khơi gợi}
(\textit{elicitation}) --- thu thập thông tin từ các bên liên quan và hiện trạng; (ii)
\emph{phân tích} (\textit{analysis}) --- chưng cất thông tin thô thành nhu cầu và xác định
mâu thuẫn, ưu tiên; (iii) \emph{đặc tả} (\textit{specification}) --- phát biểu thành yêu
cầu chức năng và phi chức năng có thể truy vết; và (iv) \emph{thẩm định}
(\textit{validation}) --- đối chiếu yêu cầu với khảo sát học thuật, công nghiệp và đánh giá
khả thi. Chương này thực hiện trọn bốn bước đó: hai bước đầu nằm ở Mục~\ref{sec:survey-enterprise}
và Mục~\ref{sec:threat-killchain}, bước đặc tả ở Mục~\ref{sec:requirements}, và bước thẩm
định trải trên các mục khảo sát và đánh giá khả thi. Kết quả --- tập yêu cầu --- sau đó
được \emph{phân tích chi tiết thành thiết kế} ở Chương~\ref{ch:design}, đúng theo sự phân
vai giữa ``xác định cần gì'' (chương này) và ``thiết kế làm thế nào'' (Chương~\ref{ch:design}).

\subsection{Môi trường và quy mô khảo sát}
\label{subsec:env-survey}

Doanh nghiệp được khảo sát vận hành một cụm Kubernetes triển khai trên hạ tầng nội bộ
(\textit{on-premise}), không dùng dịch vụ Kubernetes do nhà cung cấp đám mây quản lý (như
EKS, GKE hay AKS). Quy mô cụm ở mức lớn: hơn \textbf{một nghìn} đối tượng \texttt{Deployment}
thuộc nhiều nhóm sản phẩm, tương ứng hàng nghìn pod hoạt động đồng thời, phân bố trên
nhiều không gian tên (\textit{namespace}) theo từng đội phát triển. Mạng pod do CNI
\emph{Calico} đảm nhiệm --- một lựa chọn cho phép thực thi \texttt{NetworkPolicy} ở tầng
mạng --- song trên thực tế phần lớn không gian tên vẫn áp dụng chính sách mặc định
cho-phép-rộng, bởi đội vận hành \emph{thiếu thông tin về thực tế giao tiếp giữa các dịch
vụ} để có thể siết chặt chính sách một cách an toàn.

Về tổ chức, cụm được chia theo nhiều không gian tên ứng với từng đội sản phẩm và từng môi
trường (phát triển, kiểm thử, sản xuất). Các dịch vụ rất đa dạng về công nghệ --- nhiều
ngôn ngữ và khung lập trình cùng tồn tại --- và giao tiếp với nhau dày đặc theo mô hình
đông--đông: một yêu cầu của người dùng có thể kích hoạt một chuỗi nhiều lời gọi nội bộ
giữa các dịch vụ. Vòng đời pod ngắn (pod được tạo và hủy liên tục theo nhịp triển khai và
co giãn), kéo theo địa chỉ IP của pod thay đổi thường xuyên, khiến mọi cơ chế giám sát gắn
cứng theo địa chỉ đều nhanh chóng lỗi thời.

Về an ninh và quan trắc, doanh nghiệp đang dùng \emph{Wazuh} làm nền tảng giám sát chính.
Wazuh thu thập nhật ký, kiểm tra toàn vẹn tệp và giám sát sự kiện ở mức \emph{máy chủ}
(host) trên các nút, kết hợp với một số công cụ quản lý nhật ký tập trung và giám sát số
liệu (Prometheus/Grafana). Ở biên hệ thống có tường lửa và thiết bị kiểm soát lưu lượng
Bắc--Nam. Cách bố trí này hiệu quả với các mối đe dọa đi qua biên, nhưng có một giả định
ngầm không còn đúng trong môi trường vi dịch vụ: rằng quan sát ở tầng máy chủ và biên mạng
là \emph{đủ}. Điểm mấu chốt là toàn bộ lớp giám sát hiện có dừng ở \emph{ranh giới máy chủ
và biên mạng} --- nó hầu như không nhìn thấy lưu lượng mạng \emph{giữa các pod} bên trong
cụm, vốn là nơi diễn ra phần lớn hoạt động của một hệ vi dịch vụ quy mô lớn. Với hơn một
nghìn deployment, số lượng cặp dịch vụ có thể giao tiếp là rất lớn, và việc không có bản đồ
giao tiếp khiến cả đội an ninh lẫn đội nền tảng phải làm việc trong điều kiện thiếu thông
tin căn bản.

\subsection{Phương pháp khảo sát và thu thập yêu cầu}
\label{subsec:survey-method}

Quá trình thu thập yêu cầu được thực hiện theo cách tiếp cận quen thuộc của kỹ nghệ phần
mềm, kết hợp ba nguồn dữ liệu để giảm thiên lệch. Thứ nhất là \emph{phỏng vấn bán cấu
trúc} (\textit{semi-structured interview}) các bên liên quan; danh sách bên liên quan và
mối quan tâm chính được tóm tắt ở Bảng~\ref{tab:stakeholders}. Thứ hai là \emph{quan sát
hiện trạng}: rà soát cấu hình các công cụ giám sát đang dùng, các \texttt{NetworkPolicy}
hiện hành và bản đồ namespace của cụm. Thứ ba là \emph{phân tích sự cố quá khứ} để hiểu
những tình huống mà khả năng phát hiện hiện tại đã bỏ sót. Các câu hỏi phỏng vấn xoay
quanh bốn trục: khả năng nhìn thấy lưu lượng nội cụm, khả năng phát hiện hành vi bất
thường, quy trình ứng phó sự cố, và các ràng buộc về tài nguyên và vận hành. Cụ thể, một
số câu hỏi tiêu biểu được đặt ra gồm: ``Khi một dịch vụ bị nghi ngờ tổn hại, đội có thể
biết nó đã kết nối tới những dịch vụ nào không?''; ``Hiện có cảnh báo nào khi một pod đột
ngột tạo nhiều kết nối tới nhiều đích khác nhau không?''; ``Quy trình cách ly một pod khi
xảy ra sự cố hiện diễn ra thế nào và mất bao lâu?''; và ``Ràng buộc nào khiến việc cài
thêm tác nhân giám sát vào ứng dụng trở nên khó khăn?''.

Tổng hợp câu trả lời cho thấy một bức tranh nhất quán giữa các vai trò. Đội an ninh xác
nhận họ \emph{không} trả lời được câu hỏi về quan hệ kết nối nội cụm và \emph{không} có
cảnh báo cho hành vi dò quét bên trong; đội nền tảng nhấn mạnh rằng cài tác nhân vào hơn
một nghìn deployment đa ngôn ngữ là không khả thi về công sức bảo trì; đội vận hành lo ngại
mọi cơ chế cách ly tự động vì rủi ro chặn nhầm dịch vụ hợp lệ ở quy mô lớn; còn lập trình
viên mong muốn giải pháp ``trong suốt'' với mã ứng dụng. Những phát hiện này về sau định
hình trực tiếp các yêu cầu phi chức năng (không xâm lấn, an toàn khi phản hồi) ở
Mục~\ref{sec:requirements}.

\begin{table}[H]
  \centering
  \caption{Các bên liên quan được phỏng vấn và mối quan tâm chính.}
  \label{tab:stakeholders}
  \small
  \begin{tabularx}{\linewidth}{@{}l l X@{}}
    \toprule
    \textbf{Vai trò} & \textbf{Viết tắt} & \textbf{Mối quan tâm chính} \\
    \midrule
    Kỹ sư an ninh        & SecOps   & Phát hiện tấn công nội cụm, ứng phó sự cố, bằng chứng kiểm toán \\
    Kỹ sư nền tảng       & DevOps   & Chi phí tài nguyên, độ phức tạp vận hành, khả năng tái lập \\
    Kỹ sư vận hành       & SRE      & Tác động hiệu năng, độ ổn định, tránh báo động giả gây gián đoạn \\
    Lập trình viên       & Dev      & Không phải sửa mã ứng dụng, không phụ thuộc ngôn ngữ \\
    \bottomrule
  \end{tabularx}
\end{table}

\subsection{Hiện trạng giám sát và các điểm đau}
\label{subsec:pain-points}

Tổng hợp từ ba nguồn khảo sát, hiện trạng giám sát an ninh nội cụm bộc lộ một số điểm đau
(\textit{pain point}) nổi bật, được liệt kê ở Bảng~\ref{tab:pain-points}. Điểm đau bao
trùm và nghiêm trọng nhất là \emph{mù lưu lượng Đông--Tây}: vì Wazuh và các công cụ hiện
có chỉ quan sát ở tầng máy chủ, đội an ninh không trả lời được câu hỏi cơ bản ``dịch vụ
nào đang nói chuyện với dịch vụ nào''. Hệ quả dây chuyền là không thể phát hiện hành vi
quét cổng hay dịch chuyển ngang khi kẻ tấn công đã ở \emph{bên trong} cụm, cũng không thể
ước lượng bán kính ảnh hưởng khi một pod bị nghi ngờ chiếm quyền. Đáng chú ý, đây vừa là
bài toán an ninh vừa là bài toán vận hành: chính sự thiếu vắng bản đồ giao tiếp khiến đội
nền tảng không dám siết \texttt{NetworkPolicy} của Calico, tạo thành một vòng luẩn quẩn
giữa khả năng quan sát và khả năng kiểm soát. Hình~\ref{fig:current-gap} minh họa hiện
trạng này: lớp giám sát hiện có bao phủ được trục Bắc--Nam, nhưng toàn bộ lưu lượng
Đông--Tây giữa các pod nằm trong một \emph{vùng mù}.

\begin{figure}[H]
  \centering
  \resizebox{0.82\linewidth}{!}{% Hien trang giam sat tai doanh nghiep: Wazuh o tang host, mu luu luong pod<->pod. Chuong 1.
\begin{tikzpicture}[font=\footnotesize, >={Stealth[length=2.2mm]},
    fw/.style  ={draw, fill=red!12, align=center, minimum width=1.7cm, minimum height=1.4cm},
    pod/.style ={draw, rounded corners=2pt, fill=orange!16, align=center,
                 minimum width=1.4cm, minimum height=0.66cm},
    waz/.style ={draw, rounded corners=2pt, fill=blue!12, align=center,
                 minimum width=3.6cm, minimum height=0.6cm, font=\scriptsize},
    zone/.style={draw, dashed, rounded corners=3pt},
    ew/.style  ={<->, red!75!black, very thick, dashed},
  ]
  \node[fw] (fw) {Tường lửa\\biên\\(Bắc--Nam)};
  \node[align=center, font=\scriptsize, above=0.05cm of fw] {\textit{giám sát}\\\textit{Bắc--Nam}};

  % Node 1
  \node[pod, right=1.6cm of fw, yshift=1.55cm] (a) {svc A};
  \node[pod, right=0.55cm of a] (b) {svc B};
  \node[waz, below=0.28cm of a, xshift=0.7cm] (w1) {Wazuh agent (tầng host)};
  \node[zone, fit=(a)(b)(w1), inner sep=6pt, label={[font=\scriptsize\bfseries]above:Node 1}] (n1) {};
  % Node 2
  \node[pod, below=1.7cm of a] (c) {svc C};
  \node[pod, right=0.55cm of c] (d) {CSDL};
  \node[waz, below=0.28cm of c, xshift=0.7cm] (w2) {Wazuh agent (tầng host)};
  \node[zone, fit=(c)(d)(w2), inner sep=6pt, label={[font=\scriptsize\bfseries]below:Node 2}] (n2) {};

  \node[draw, thick, rounded corners=4pt, fit=(n1)(n2), inner sep=11pt,
        label={[font=\bfseries]above:Cụm Kubernetes (CNI Calico)}] (cl) {};

  \draw[->] (fw.east) -- (a.west);
  % Luu luong Dong-Tay KHONG duoc giam sat
  \draw[ew] (a) -- (b);
  \draw[ew] (b.south) to[bend left=14] (d.north);
  \draw[ew] (c) -- (d);
  \draw[ew] (a.south) to[bend right=10] (c.north);

  \node[draw=red!70!black, fill=red!6, rounded corners, align=left, font=\scriptsize,
        text width=3.0cm, right=1.1cm of cl] (blind)
    {\textbf{Vùng mù:} lưu lượng pod$\leftrightarrow$pod (Đông--Tây) \emph{không} được giám sát;
     Wazuh chỉ thấy tầng host.};
  \draw[->, red!70!black] (blind.west) -- (cl.east);
\end{tikzpicture}
}
  \caption{Hiện trạng giám sát và vùng mù Đông--Tây tại doanh nghiệp.}
  \label{fig:current-gap}
\end{figure}

\begin{table}[H]
  \centering
  \caption{Các điểm đau ghi nhận từ khảo sát hiện trạng.}
  \label{tab:pain-points}
  \small
  \begin{tabularx}{\linewidth}{@{}l l X@{}}
    \toprule
    \textbf{Mã} & \textbf{Điểm đau} & \textbf{Mô tả} \\
    \midrule
    P1 & Mù lưu lượng Đông--Tây   & Không thấy giao tiếp pod-tới-pod; chỉ giám sát ở tầng host (Wazuh) và biên mạng \\
    P2 & Thiếu bản đồ phụ thuộc   & Không biết dịch vụ nào gọi dịch vụ nào nên không siết được \texttt{NetworkPolicy} \\
    P3 & Tác nhân theo pod bất khả thi & Cài agent vào hơn 1000 deployment là tốn kém, phụ thuộc ngôn ngữ, khó bảo trì \\
    P4 & Ứng phó thủ công, chậm   & Khi nghi ngờ sự cố, việc cách ly pod làm thủ công, không có quy trình tự động có kiểm soát \\
    P5 & Thiếu bằng chứng kiểm toán & Không có dấu vết ``ai kết nối tới ai'' trong cụm phục vụ điều tra và tuân thủ \\
    \bottomrule
  \end{tabularx}
\end{table}

\subsection{Tổng hợp nhu cầu}
\label{subsec:needs}

Từ các điểm đau, quá trình khảo sát chưng cất ra một tập \emph{nhu cầu}
(Bảng~\ref{tab:needs}) --- là phát biểu ``cần gì'' ở mức nghiệp vụ, độc lập với giải pháp
kỹ thuật. Tập nhu cầu này định hình \emph{miền bài toán} (domain) thực sự mà đồ án phải
giải quyết, và sẽ được chuyển hóa thành các yêu cầu chức năng và phi chức năng cụ thể ở
Mục~\ref{sec:requirements}, rồi tiếp tục được phân tích chi tiết thành thiết kế ở
Chương~\ref{ch:design}.

\begin{table}[H]
  \centering
  \caption{Tổng hợp nhu cầu chưng cất từ các điểm đau.}
  \label{tab:needs}
  \small
  \begin{tabularx}{\linewidth}{@{}l X l@{}}
    \toprule
    \textbf{Mã} & \textbf{Nhu cầu} & \textbf{Đáp ứng điểm đau} \\
    \midrule
    N1 & Quan sát lưu lượng kết nối toàn cụm với chi phí thấp, \emph{không} cài tác nhân vào từng pod & P1, P3 \\
    N2 & Phát hiện trực tuyến hành vi quét cổng và dịch chuyển ngang & P1 \\
    N3 & Dựng bản đồ giao tiếp giữa các dịch vụ và ước lượng bán kính ảnh hưởng & P2 \\
    N4 & Ứng phó cách ly tự động \emph{có người duyệt}, tận dụng \texttt{NetworkPolicy} của Calico sẵn có & P2, P4 \\
    N5 & Lưu vết kiểm toán đầy đủ phục vụ điều tra và tuân thủ & P5 \\
    \bottomrule
  \end{tabularx}
\end{table}

\section{Mô hình mối đe dọa: chuỗi tấn công nội cụm}
\label{sec:threat-killchain}

Để hiểu vì sao các nhu cầu trên lại tập trung vào giai đoạn ``bên trong'' cụm, cần đặt
chúng vào một mô hình mối đe dọa. Một cuộc tấn công vào hệ vi dịch vụ thường diễn ra theo
một chuỗi giai đoạn (\textit{kill-chain}) như Hình~\ref{fig:kill-chain}: sau khi
\emph{xâm nhập ban đầu} qua một pod lộ lỗ hổng (giai đoạn~1), kẻ tấn công \emph{trinh sát}
bằng cách quét cổng các pod và dịch vụ lân cận (giai đoạn~2), rồi \emph{dịch chuyển ngang}
từ pod này sang pod khác (giai đoạn~3) để tiến tới \emph{mục tiêu giá trị cao} như cơ sở
dữ liệu hay kho bí mật (giai đoạn~4), và cuối cùng \emph{tác động} bằng việc rò rỉ hay phá
hoại dữ liệu (giai đoạn~5).

\begin{figure}[H]
  \centering
  \resizebox{0.96\linewidth}{!}{% Chuoi tan cong (kill-chain) noi cum K8s; danh dau vung do an phat hien/ngan chan. Chuong 1.
\begin{tikzpicture}[font=\footnotesize, >={Stealth[length=2.6mm]},
    st/.style ={draw, rounded corners=3pt, align=center, minimum width=2.55cm,
                minimum height=1.25cm, fill=gray!8},
    hot/.style={draw, rounded corners=3pt, align=center, minimum width=2.55cm,
                minimum height=1.25cm, fill=red!12, draw=red!60!black, line width=0.9pt},
    flow/.style={->, very thick, gray!55},
  ]
  \node[st]  (s1) at (0,0)     {\textbf{1. Xâm nhập}\\ban đầu\\\scriptsize(pod lộ lỗ hổng)};
  \node[hot] (s2) at (3.0,0)   {\textbf{2. Trinh sát}\\quét cổng\\\scriptsize(nội cụm)};
  \node[hot] (s3) at (6.0,0)   {\textbf{3. Dịch chuyển}\\\textbf{ngang}\\\scriptsize(pod $\rightarrow$ pod)};
  \node[st]  (s4) at (9.0,0)   {\textbf{4. Tiếp cận}\\mục tiêu\\\scriptsize(CSDL, secret)};
  \node[st]  (s5) at (12.0,0)  {\textbf{5. Tác động}\\\scriptsize(rò rỉ /\\phá hoại)};

  \draw[flow] (s1)--(s2); \draw[flow] (s2)--(s3); \draw[flow] (s3)--(s4); \draw[flow] (s4)--(s5);

  % Vung phat hien cua do an = giai doan 2-3
  \draw[decorate, decoration={brace, amplitude=7pt, mirror}]
    ($(s2.south west)+(0,-0.2)$) -- ($(s3.south east)+(0,-0.2)$)
    node[midway, below=10pt, align=center, font=\footnotesize\bfseries, text=red!60!black]
    {Vùng K8sSecDaP phát hiện sớm \& ngăn chặn\\(tín hiệu tần suất $+$ cấu trúc đồ thị)};

  % Ghi chu an ninh chu vi chi chan duoc 2 dau
  \draw[decorate, decoration={brace, amplitude=6pt}]
    ($(s1.north west)+(0,0.2)$) -- ($(s1.north east)+(0,0.2)$)
    node[midway, above=8pt, font=\scriptsize\itshape] {an ninh chu vi};
  \draw[decorate, decoration={brace, amplitude=6pt}]
    ($(s5.north west)+(0,0.2)$) -- ($(s5.north east)+(0,0.2)$)
    node[midway, above=8pt, font=\scriptsize\itshape] {DLP / lối ra};
\end{tikzpicture}
}
  \caption{Chuỗi tấn công (kill-chain) nội cụm và vùng phát hiện của đồ án.}
  \label{fig:kill-chain}
\end{figure}

Điểm mấu chốt là an ninh chu vi truyền thống chỉ kiểm soát được hai đầu của chuỗi: ngăn
xâm nhập ban đầu ở biên (giai đoạn~1) và chặn rò rỉ ra ngoài (giai đoạn~5). Toàn bộ phần
\emph{giữa} --- trinh sát và dịch chuyển ngang --- diễn ra hoàn toàn trong lưu lượng
Đông--Tây và vô hình trước lớp giám sát hiện tại của doanh nghiệp. Đây chính xác là khoảng
mà đồ án nhắm tới: phát hiện \emph{sớm} ở giai đoạn~2--3, khi kẻ tấn công vừa mới bắt đầu
dò quét và lan tỏa, để ngăn chặn trước khi chúng chạm tới mục tiêu. Hai giai đoạn này để
lại hai loại tín hiệu bổ trợ đã nêu ở Mục~\ref{subsec:threat-defs} (tần suất kết nối và
cấu trúc đồ thị giao tiếp), và đó là nền tảng cho cách tiếp cận phát hiện của hệ thống.

\subsection{Kịch bản minh họa trong bối cảnh doanh nghiệp}

Để cụ thể hóa, hãy xét một kịch bản đặt trong đúng môi trường đã khảo sát. Một pod chạy
dịch vụ web đối mặt người dùng tồn tại lỗ hổng và bị kẻ tấn công kiểm soát (giai đoạn~1).
Từ chỗ đứng này, kẻ tấn công không tấn công thẳng ra biên --- nơi vẫn được tường lửa canh
gác --- mà quay vào trong: chạy một công cụ quét cổng nhắm tới dải địa chỉ pod của các
namespace lân cận để tìm dịch vụ đang mở (giai đoạn~2). Với lớp giám sát hiện tại, hàng
trăm lượt kết nối dò quét này hòa lẫn vào lưu lượng Đông--Tây và \emph{không} sinh ra cảnh
báo nào, bởi Wazuh chỉ thấy sự kiện ở tầng host. Phát hiện được mẫu đang mở, kẻ tấn công
lần lượt kết nối sang một pod xử lý nội bộ rồi sang pod cơ sở dữ liệu (giai đoạn~3--4),
hình thành một chuỗi liên lạc pod-tới-pod bất thường mà không công cụ nào trong cụm ghi
nhận.

Đặt cạnh nhau, kịch bản này làm lộ rõ hai cơ hội phát hiện mà hệ thống đề xuất khai thác.
Thứ nhất, ở giai đoạn trinh sát, pod tấn công tạo ra một \emph{số lượng kết nối tăng vọt}
trong thời gian ngắn --- một dị thường về tần suất mà bộ đếm Count-Min Sketch nắm bắt được
ngay khi sự kiện vừa sinh ra. Thứ hai, ở giai đoạn dịch chuyển ngang, chuỗi pod bị chiếm
tạo ra một \emph{cấu trúc liên thông bất thường} trên đồ thị giao tiếp mà giải thuật đồ thị
phát hiện được, đồng thời cho phép ước lượng bán kính ảnh hưởng để ưu tiên xử lý. Chính
việc một sự kiện đơn lẻ đã đủ để kích hoạt tín hiệu tần suất, còn cấu trúc đồ thị bổ sung
bức tranh quan hệ, là lý do hệ thống có thể can thiệp \emph{trước} khi kẻ tấn công chạm
tới mục tiêu --- điều mà lớp giám sát theo lô, hướng host hiện tại không làm được.

\section{Khảo sát hiện trạng và khoảng trống nghiên cứu}

\subsection{Bằng chứng từ ba bài tổng quan có hệ thống gần đây}

Nhu cầu thực tế ghi nhận ở doanh nghiệp không phải là một trường hợp cá biệt. Để đặt nó
vào bức tranh chung và bảo đảm đề tài không xuất phát từ giả định chủ quan, đồ án đối
chiếu với \textbf{ba bài khảo sát có hệ thống} (\textit{systematic literature review}, SLR)
gần nhất về an toàn thông tin của hệ vi dịch vụ trên Kubernetes. SLR là phương pháp tổng
quan tài liệu theo quy trình minh bạch (xác định câu hỏi, sàng lọc theo tiêu chí, tổng hợp
có hệ thống), nên kết luận của chúng có độ tin cậy cao hơn các bài tổng quan thông thường.
Ba công trình dưới đây bổ sung cho nhau cả về quy mô khảo sát lẫn góc nhìn:

\textbf{Hutasuhut và cộng sự (2024)} \cite{hutasuhut2024gap} sàng lọc 487 công bố trong
giai đoạn sau năm 2020 và chọn lại 87 bài bằng phương pháp snowball. Nhóm tác giả kết luận
rằng vẫn thiếu vắng đáng kể các chuẩn bảo mật riêng và các mô hình toàn diện được thiết kế
\emph{đặc biệt} cho vi dịch vụ, đồng thời tồn tại một khoảng cách kiến thức rõ rệt giữa giới
học thuật và giới thực hành công nghiệp. Khoảng cách này trùng khớp với điều quan sát được ở
doanh nghiệp: các đội vận hành biết mình thiếu khả năng quan sát nội cụm nhưng không có một
mô hình tham chiếu sẵn để áp dụng.

\textbf{Berardi và cộng sự (2022)} \cite{berardi2022microservice} khảo sát 290 công bố và
xác định bốn khoảng trống cốt lõi: thiếu hướng dẫn áp dụng \textit{security-by-design} cho
kiến trúc vi dịch vụ; thiếu mô hình mối đe dọa chuyên biệt; thiếu hướng dẫn xử lý bề mặt
tấn công sinh ra từ sự đa dạng công nghệ; và các vấn đề an toàn khi di trú từ hệ đơn khối
sang vi dịch vụ. Đáng chú ý, ba trong bốn khoảng trống này liên quan trực tiếp tới mô hình
mối đe dọa và bề mặt tấn công --- đúng phần mà đồ án tập trung làm rõ qua chuỗi kill-chain
và hai mẫu hành vi quét cổng, dịch chuyển ngang.

\textbf{Pereira-Vale và cộng sự (2021)} \cite{pereiravale2021security} thực hiện một bài
tổng quan đa nguồn (\textit{multivocal review}) kết hợp 36 công bố học thuật với 34 nguồn
xám (báo cáo công nghiệp, tài liệu kỹ thuật của các nhà cung cấp lớn). Kết quả cho thấy
một sự lệch pha hai chiều: nhiều giải pháp được công nghiệp triển khai chưa có cơ sở lý
thuyết được kiểm chứng, trong khi nhiều phân tích lý thuyết lại chưa được hiện thực hóa
trong hệ thống thực. Chính sự lệch pha này củng cố cho định hướng của đồ án --- gắn chặt
mỗi giải thuật cốt lõi với cơ sở lý thuyết \emph{và} một hiện thực đo được trên cụm thật.

\subsection{Đối chiếu với báo cáo công nghiệp gần đây}

Để đối chiếu kết luận học thuật với thực tế triển khai, đồ án tham chiếu hai báo cáo công nghiệp gần nhất:

\begin{itemize}
  \item \textit{The State of Kubernetes Security 2024} của Red Hat \cite{redhat2024k8ssec}, dựa trên khảo sát 600 chuyên gia DevOps và an toàn thông tin: \textbf{67\%} tổ chức đã trì hoãn hoặc làm chậm phát triển ứng dụng do lo ngại an ninh; \textbf{45\%} gặp sự cố ở giai đoạn \textit{runtime} trong 12 tháng gần nhất; \textbf{30\%} bị phạt và \textbf{46\%} chịu mất doanh thu hoặc khách hàng sau sự cố.
  \item \textit{Kubernetes Security Report 2025} của Wiz \cite{wiz2025k8ssec} bổ sung số liệu thời gian: một cluster K8s mới tạo trên dịch vụ managed (AKS, EKS) bị tấn công thử lần đầu trong vòng 18--28 phút sau khi sẵn sàng, cho thấy ``đầu vào đối kháng'' không phải tình huống hiếm gặp mà là điều kiện vận hành mặc định.
\end{itemize}

Bảng~\ref{tab:survey-data} tổng hợp các con số tiêu biểu từ các nguồn khảo sát trên,
cho thấy an ninh Kubernetes vừa là vấn đề phổ biến vừa có hậu quả kinh doanh rõ rệt --- một
luận cứ định lượng củng cố cho nhu cầu thực tế đã ghi nhận ở
Mục~\ref{sec:survey-enterprise}.

\begin{table}[H]
  \centering
  \caption{Một số số liệu tiêu biểu về an ninh Kubernetes từ các khảo sát gần đây.}
  \label{tab:survey-data}
  \small
  \begin{tabularx}{\linewidth}{@{}l l X@{}}
    \toprule
    \textbf{Nguồn} & \textbf{Số liệu} & \textbf{Ý nghĩa} \\
    \midrule
    Hutasuhut 2024 \cite{hutasuhut2024gap} & 487 $\rightarrow$ 87 bài & Thiếu chuẩn bảo mật riêng và mô hình toàn diện cho vi dịch vụ \\
    Berardi 2022 \cite{berardi2022microservice} & 290 bài, 4 khoảng trống & Thiếu security-by-design và mô hình mối đe dọa chuyên biệt \\
    Red Hat 2024 \cite{redhat2024k8ssec} & 67\% / 45\% / 30\% / 46\% & Trì hoãn phát hành / sự cố runtime / bị phạt / mất doanh thu \\
    Wiz 2025 \cite{wiz2025k8ssec} & 18--28 phút & Thời gian tới lần dò quét đầu tiên của một cụm mới \\
    \bottomrule
  \end{tabularx}
\end{table}

\subsection{Khoảng trống cụ thể mà đồ án nhắm tới}

Khi đặt ba nguồn khảo sát cạnh nhau --- nhu cầu thực tế của doanh nghiệp, các bài tổng
quan học thuật và các công cụ công nghiệp --- chúng hội tụ về cùng một khoảng trống. Từ
phía \emph{doanh nghiệp}, đó là sự mù lưu lượng Đông--Tây và vòng luẩn quẩn quan
sát -- kiểm soát. Từ phía \emph{học thuật}, đó là sự thiếu vắng mô hình mối đe dọa chuyên
biệt và kiến trúc tham chiếu cho vi dịch vụ. Từ phía \emph{công cụ}, đó là việc mỗi công cụ
chỉ giải quyết một mảnh và hiếm khi lấy cấu trúc dữ liệu luồng làm trọng tâm phát hiện.

Tổng hợp lại, đồ án xác định một \textit{trục} khoảng trống cụ thể: thiếu một \textbf{kiến
trúc tham chiếu hoàn chỉnh} kết nối từ tầng nhân (thu thập sự kiện mạng), qua tầng cấu trúc
dữ liệu (giải thuật tần suất, đồ thị), tới tầng dịch vụ (phát hiện sự cố, phản ứng) và giao
diện vận hành (dashboard và console). Hầu hết các giải pháp công cụ hiện có (như Cilium
Hubble, Falco, Tetragon) tập trung vào một mảnh và không cho thấy cách các mảnh đó được
\textit{ráp lại} thành một quy trình giám sát Zero-Trust khép kín. Đồ án nhằm bù đắp đúng
khoảng trống này bằng cách xây dựng và đánh giá một hệ tham chiếu đầu cuối, được thiết kế
để các thành phần lý thuyết và kỹ thuật có thể được tách rời và tái sử dụng --- vừa giải
quyết bài toán cụ thể của doanh nghiệp, vừa đóng góp một mẫu hình có thể khái quát hóa.

\section{Mục tiêu, đối tượng và phạm vi}

\subsection{Mục tiêu nghiên cứu}

\begin{enumerate}
    \item Thu thập sự kiện kết nối TCP ở \textbf{tầng nhân Linux bằng eBPF} (\texttt{kprobe tcp\_v4\_connect}, \texttt{BPF\_MAP\_TYPE\_LPM\_TRIE}, \texttt{BPF\_MAP\_TYPE\_RINGBUF}) trên mọi nút với chi phí thấp.
    \item Hình thức hóa và giải quyết các bài toán cấu trúc dữ liệu luồng cốt lõi: ước lượng tần suất (Count-Min Sketch), phát hiện thành phần liên thông mạnh (Tarjan), khả-đạt-k-bước (BFS), phân loại địa chỉ IP theo tiền tố (LPM Trie).
    \item Thiết kế kiến trúc phần mềm phân tầng, cấu hình hóa theo YAML, dùng mẫu thiết kế Strategy~+~Registry cho phép đổi giải thuật không cần sửa mã nguồn.
    \item Xây dựng hệ thống SOC nội bộ (NATS, PostgreSQL, dịch vụ incident, bộ sinh NetworkPolicy theo cơ chế \textit{generate-for-review}) thực hiện vòng phát hiện$\rightarrow$phản hồi có người duyệt.
    \item \textbf{Đánh giá định lượng} trên cluster K8s thật: tỉ lệ phát hiện, tỉ lệ báo động giả, chi phí eBPF, độ trễ phát hiện$\rightarrow$phản hồi.
\end{enumerate}

\subsection{Đối tượng và phạm vi}

\textbf{Đối tượng nghiên cứu:} hành vi quét cổng và dịch chuyển ngang trong mạng nội cụm K8s; các cấu trúc dữ liệu luồng và đồ thị phục vụ phát hiện; cơ chế thực thi chính sách mạng (NetworkPolicy).

\noindent\textbf{Phạm vi:}
\begin{itemize}
    \item Kiến trúc Zero-Trust cho môi trường K8s on-premise (không bao gồm cloud-managed như EKS/GKE).
    \item Phát hiện các mẫu tấn công tiêu biểu: quét cổng (\textit{port scan}), dịch chuyển ngang (\textit{lateral movement}) liên nút.
    \item Toàn bộ thành phần do sinh viên tự cài đặt, ngoại trừ hạ tầng quan trắc và message bus dùng bản chuẩn thượng nguồn.
\end{itemize}

\noindent\textbf{Giới hạn:}
\begin{itemize}
    \item Phản ứng ở chế độ \emph{generate-for-review}: hệ thống tạo bản nháp NetworkPolicy, người vận hành phê duyệt; không tự động áp đặt.
    \item Chỉ hỗ trợ IPv4 (\texttt{tcp\_v4\_connect}); đánh giá ở quy mô lớn là hướng phát triển.
\end{itemize}

Bảng~\ref{tab:scope} tóm tắt ranh giới phạm vi để tránh hiểu nhầm về kỳ vọng đối với hệ
thống.

\begin{table}[H]
  \centering
  \caption{Ranh giới phạm vi của đồ án.}
  \label{tab:scope}
  \small
  \begin{tabularx}{\linewidth}{@{}X X@{}}
    \toprule
    \textbf{Trong phạm vi} & \textbf{Ngoài phạm vi} \\
    \midrule
    Thu thập kết nối TCP/IPv4 ở tầng nhân & Lưu lượng UDP, ICMP, IPv6 \\
    Phát hiện quét cổng và dịch chuyển ngang & Phát hiện mã độc, khai thác lỗ hổng tầng ứng dụng \\
    Cụm Kubernetes on-premise & Dịch vụ K8s do đám mây quản lý (EKS/GKE/AKS) \\
    Phản hồi cách ly có người duyệt & Tự động cách ly không người duyệt \\
    Đánh giá trên cụm ba nút & Đánh giá tải ở quy mô hàng nghìn nút \\
    \bottomrule
  \end{tabularx}
\end{table}

\section{Khảo sát giải pháp và công cụ liên quan}
\label{sec:related-work}

Để định vị đề tài, mục này khảo sát hai nhóm giải pháp: an ninh \emph{truyền thống} và an
ninh \emph{cloud-native}, rồi đối chiếu các công cụ tiêu biểu đang được công nghiệp dùng.
Phần so sánh hiệu năng định lượng chi tiết với các công cụ này được đặt ở
Chương~\ref{ch:eval}; ở đây tập trung vào khác biệt về \emph{cách tiếp cận}.

\subsection{An ninh truyền thống so với an ninh cloud-native}

An ninh truyền thống dựa trên ba trụ: tường lửa/IDS đặt ở biên mạng, tác nhân giám sát
cài trên từng máy chủ (host-agent, ví dụ Wazuh trong môi trường khảo sát), và hệ SIEM
phân tích nhật ký theo lô. Mô hình này phù hợp với kiến trúc một khối nhưng bộc lộ giới
hạn rõ rệt trong môi trường vi dịch vụ: nó không quan sát được lưu lượng Đông--Tây nội cụm,
không bám theo vòng đời pod ngắn, và phản ứng theo lô nên trễ. Bảng~\ref{tab:trad-vs-cloud}
đối chiếu hai cách tiếp cận trên các trục then chốt.

\begin{table}[H]
  \centering
  \caption{So sánh an ninh truyền thống và an ninh cloud-native (tầng nhân).}
  \label{tab:trad-vs-cloud}
  \small
  \begin{tabularx}{\linewidth}{@{}l X X@{}}
    \toprule
    \textbf{Trục} & \textbf{Truyền thống} & \textbf{Cloud-native (đồ án)} \\
    \midrule
    Vị trí quan sát   & Biên mạng / máy chủ (Bắc--Nam)        & Tầng nhân mỗi nút (Đông--Tây) \\
    Cách thu thập     & Host-agent, mirror gói, theo lô       & eBPF không xâm lấn, một \texttt{DaemonSet}/nút \\
    Độ phủ nội cụm    & Hầu như không thấy pod-tới-pod        & Toàn bộ kết nối pod-tới-pod \\
    Phụ thuộc ứng dụng& Có (agent theo ngôn ngữ)              & Không (ở tầng nhân) \\
    Mô hình vận hành  & Cấu hình thủ công, dễ trôi dạt        & Khai báo GitOps, ảnh bất biến, tái lập \\
    Vòng phản hồi     & Phân tích lô, xử lý thủ công          & Phát hiện trực tuyến, cách ly có người duyệt \\
    \bottomrule
  \end{tabularx}
\end{table}

\subsection{Các công cụ an ninh cloud-native tiêu biểu}

Trong hệ sinh thái cloud-native đã xuất hiện nhiều công cụ khai thác eBPF hoặc BPF-LSM.
\emph{Falco} \cite{falco2024} phát hiện bất thường theo luật trên dòng lời gọi hệ thống;
\emph{Tetragon} \cite{tetragon2024} cung cấp quan sát và cưỡng chế ở không gian nhân qua
LSM; \emph{Cilium} \cite{cilium2024} là một CNI dựa trên eBPF với khả năng quan sát luồng
(Hubble) và \texttt{NetworkPolicy} L3--L7; \emph{KubeArmor} \cite{kubearmor2024} cưỡng chế
hành vi tiến trình/tệp/mạng bằng BPF-LSM. Các công cụ NIDS truyền thống như Snort
\cite{roesch1999snort} và Bro/Zeek \cite{paxson1999bro} phân tích gói ở mức mạng nhưng
không gắn với ngữ cảnh điều phối của Kubernetes. Bảng~\ref{tab:tool-compare} đối chiếu các
công cụ này với đồ án.

\begin{table}[H]
  \centering
  \caption{Đối chiếu một số công cụ liên quan với đồ án.}
  \label{tab:tool-compare}
  \small
  \begin{tabularx}{\linewidth}{@{}l X c c@{}}
    \toprule
    \textbf{Công cụ} & \textbf{Trọng tâm} & \textbf{eBPF} & \textbf{CTDL luồng} \\
    \midrule
    Falco          & Phát hiện theo luật trên syscall        & Có   & Không \\
    Tetragon       & Quan sát + cưỡng chế qua LSM            & Có   & Không \\
    Cilium/Hubble  & CNI + quan sát luồng + policy L3--L7    & Có   & Không \\
    KubeArmor      & Cưỡng chế hành vi qua BPF-LSM           & Có   & Không \\
    Wazuh          & SIEM/host-IDS theo nhật ký             & Không& Không \\
    Snort, Bro/Zeek& NIDS phân tích gói ở mức mạng          & Không& Không \\
    \textbf{Đồ án} & \textbf{Phát hiện quét cổng/dịch chuyển ngang} & \textbf{Có} & \textbf{Có (CMS, đồ thị)} \\
    \bottomrule
  \end{tabularx}
\end{table}

Điểm chung của các công cụ trên là khai thác eBPF/LSM để quan sát hoặc cưỡng chế ở tầng
nhân. Điểm khác biệt --- cũng là khoảng trống đồ án lấp --- là rất ít công cụ đặt
\emph{cấu trúc dữ liệu luồng} (Count-Min Sketch) và \emph{giải thuật đồ thị} (Tarjan, BFS)
làm trọng tâm phát hiện quét cổng và dịch chuyển ngang, đồng thời ráp các mảnh thành một
quy trình phát hiện$\rightarrow$phản hồi khép kín có người duyệt.

\section{Phân tích yêu cầu, đánh giá khả thi và đóng góp}
\label{sec:requirements}

Từ tập nhu cầu nghiệp vụ ở Mục~\ref{subsec:needs}, mục này chưng cất thành các \emph{yêu
cầu} cụ thể --- bước chuyển từ ``cần gì'' sang ``hệ thống phải làm gì''. Đây là đầu ra của
giai đoạn phân tích yêu cầu ở mức tổng quan; việc phân tích chi tiết thành ca sử dụng, sơ
đồ và thiết kế được tiếp nối ở Chương~\ref{ch:design}. Hình~\ref{fig:re-process} tóm tắt
mạch truy vết xuyên suốt: từ khảo sát nhu cầu, qua đặc tả yêu cầu, tới thiết kế, cài đặt và
đánh giá --- trong đó chương này phụ trách phần ``cần gì''.

\begin{figure}[H]
  \centering
  \resizebox{0.96\linewidth}{!}{% Quy trinh ky nghe yeu cau va truy vet sang cac chuong sau. Chuong 1.
\begin{tikzpicture}[font=\footnotesize, >={Stealth[length=2.4mm]}, node distance=0.55cm,
    s1/.style={draw, rounded corners=3pt, align=center, fill=blue!10,
               minimum width=2.5cm, minimum height=1.05cm},
    s2/.style={draw, rounded corners=3pt, align=center, fill=green!12,
               minimum width=2.5cm, minimum height=1.05cm},
    flow/.style={->, very thick},
  ]
  \node[s1] (a) {Khảo sát\\nhu cầu\\\scriptsize(\S\ref{sec:survey-enterprise})};
  \node[s1, right=of a] (b) {Nhu cầu\\N1--N5\\\scriptsize(\S\ref{subsec:needs})};
  \node[s1, right=of b] (c) {Yêu cầu\\FR / NFR\\\scriptsize(\S\ref{sec:requirements})};
  \node[s2, right=of c] (d) {Thiết kế\\\scriptsize(Ch.~\ref{ch:design})};
  \node[s2, right=of d] (e) {Cài đặt\\\scriptsize(Ch.~\ref{ch:impl})};
  \node[s2, right=of e] (f) {Đánh giá\\\scriptsize(Ch.~\ref{ch:eval})};

  \foreach \x/\y in {a/b,b/c,c/d,d/e,e/f}{ \draw[flow] (\x)--(\y); }

  \draw[decorate, decoration={brace, amplitude=6pt}]
    ($(a.north west)+(0,0.2)$) -- ($(c.north east)+(0,0.2)$)
    node[midway, above=8pt, font=\footnotesize\bfseries] {Chương~\ref{ch:overview} (xác định ``cần gì'')};
  \draw[decorate, decoration={brace, amplitude=6pt, mirror}]
    ($(d.south west)+(0,-0.2)$) -- ($(f.south east)+(0,-0.2)$)
    node[midway, below=8pt, font=\footnotesize\itshape] {``làm thế nào''};
\end{tikzpicture}
}
  \caption{Quy trình kỹ nghệ yêu cầu và mạch truy vết từ nhu cầu tới các chương sau.}
  \label{fig:re-process}
\end{figure}

\subsection{Yêu cầu chức năng}

Bảng~\ref{tab:fr} liệt kê các yêu cầu chức năng (\textit{functional requirement}), mỗi
yêu cầu truy vết được về nhu cầu đã nêu.

\begin{table}[H]
  \centering
  \caption{Yêu cầu chức năng của hệ thống.}
  \label{tab:fr}
  \small
  \begin{tabularx}{\linewidth}{@{}l X l@{}}
    \toprule
    \textbf{Mã} & \textbf{Yêu cầu chức năng} & \textbf{Nhu cầu} \\
    \midrule
    FR-01 & Thu thập sự kiện kết nối TCP toàn cụm ở tầng nhân (eBPF)        & N1 \\
    FR-02 & Phân loại địa chỉ theo vùng mạng (pod/dịch vụ/nút/ngoài)        & N2 \\
    FR-03 & Phát hiện quét cổng theo tần suất kết nối (Count-Min Sketch)    & N2 \\
    FR-04 & Dựng đồ thị giao tiếp; phát hiện liên thông mạnh và bán kính ảnh hưởng & N3 \\
    FR-05 & Tạo và làm giàu sự cố (ánh xạ địa chỉ nguồn $\rightarrow$ pod)  & N3, N5 \\
    FR-06 & Sinh bản nháp \texttt{NetworkPolicy} cách ly (egress-deny)      & N4 \\
    FR-07 & Quy trình phê duyệt có người duyệt trước khi áp chính sách      & N4 \\
    FR-08 & Áp chính sách qua API Kubernetes và xác minh hiệu lực           & N4 \\
    FR-09 & Ghi nhật ký kiểm toán mọi sự cố và hành động                    & N5 \\
    FR-10 & Bảng điều khiển và console vận hành cho người dùng              & N3, N5 \\
    \bottomrule
  \end{tabularx}
\end{table}

\subsection{Yêu cầu phi chức năng}

Bảng~\ref{tab:nfr} liệt kê các yêu cầu phi chức năng (\textit{non-functional requirement})
--- những ràng buộc về chất lượng quyết định tính khả dụng của giải pháp ở quy mô doanh
nghiệp.

\begin{table}[H]
  \centering
  \caption{Yêu cầu phi chức năng của hệ thống.}
  \label{tab:nfr}
  \small
  \begin{tabularx}{\linewidth}{@{}l X@{}}
    \toprule
    \textbf{Mã} & \textbf{Yêu cầu phi chức năng} \\
    \midrule
    NFR-01 & Chi phí thấp: phần thu thập ở tầng nhân chiếm phần nhỏ một lõi CPU mỗi nút \\
    NFR-02 & Khả mở rộng: hoạt động được trên cụm hàng nghìn pod / nhiều nút \\
    NFR-03 & Không xâm lấn: không cài tác nhân vào ứng dụng, độc lập ngôn ngữ \\
    NFR-04 & Tái lập được: triển khai khai báo bằng GitOps với ảnh bất biến \\
    NFR-05 & Kiểm toán được: mọi hành động phản hồi đều để lại dấu vết \\
    NFR-06 & Tự chủ nội bộ: chạy trọn vẹn trên hạ tầng on-premise, không phụ thuộc đám mây \\
    NFR-07 & An toàn khi phản hồi: cách ly chỉ được áp sau khi có người phê duyệt \\
    \bottomrule
  \end{tabularx}
\end{table}

\subsection{Đánh giá tính khả thi}
\label{sec:feasibility}

Trước khi cam kết thiết kế, cần đánh giá tính khả thi của giải pháp đề xuất trên ba khía
cạnh. Về \emph{kỹ thuật}: các nút của cụm chạy nhân Linux hiện đại có sẵn thông tin BTF,
đủ điều kiện để nạp chương trình eBPF CO-RE; CNI Calico đã hiện diện sẵn nên cơ chế cưỡng
chế \texttt{NetworkPolicy} cho vòng phản hồi là sẵn dùng mà không phải thay đổi hạ tầng
mạng; chi phí quan sát ở tầng nhân được đo là rất nhỏ (Chương~\ref{ch:eval}), phù hợp với
ràng buộc NFR-01. Về \emph{kinh tế}: toàn bộ thành phần là mã nguồn mở, không phát sinh
chi phí giấy phép theo pod --- một lợi thế quyết định khi quy mô vượt một nghìn deployment,
nơi mô hình tác nhân-theo-pod trở nên đắt đỏ. Về \emph{vận hành}: mô hình \texttt{DaemonSet}
cho phép triển khai một lần trên mỗi nút và tự lan theo cụm; vòng phản hồi có người duyệt
tránh rủi ro tự động cách ly nhầm ở quy mô lớn, đáp ứng đúng lo ngại của đội SRE đã ghi
nhận khi khảo sát. Cả ba khía cạnh đều cho thấy giải pháp khả thi trong đúng bối cảnh
doanh nghiệp đã khảo sát.

Tính khả thi đi kèm một số giả định và rủi ro cần được nhận diện sớm
(Bảng~\ref{tab:risk}); việc nêu rõ chúng giúp khoanh vùng đúng những điều kiện mà kết luận
của đồ án còn đúng và định hướng cho phần giới hạn cùng hướng phát triển ở
Chương~\ref{ch:eval}.

\begin{table}[H]
  \centering
  \caption{Giả định, rủi ro và biện pháp giảm thiểu.}
  \label{tab:risk}
  \small
  \begin{tabularx}{\linewidth}{@{}l X X@{}}
    \toprule
    \textbf{Mã} & \textbf{Giả định / rủi ro} & \textbf{Biện pháp giảm thiểu} \\
    \midrule
    R1 & Nhân nút có BTF cho CO-RE        & Kiểm tra điều kiện nhân; sinh \texttt{vmlinux.h} lúc xây ảnh \\
    R2 & Quét gửi SYN thô né \texttt{connect()} & Ghi nhận giới hạn; bổ sung điểm móc khác là hướng phát triển \\
    R3 & Client hợp lệ tạo nhiều kết nối $\rightarrow$ báo giả & Đếm theo cửa sổ + người duyệt trước khi cách ly \\
    R4 & Cách ly nhầm gây gián đoạn dịch vụ & Cơ chế \emph{generate-for-review}, có người phê duyệt \\
    R5 & Đồ thị giao tiếp lớn dần theo thời gian & Chạy Tarjan/BFS theo chu kỳ, độ phức tạp tuyến tính \\
    \bottomrule
  \end{tabularx}
\end{table}

\subsection{Đóng góp của đồ án}
\label{sec:contributions}

Trên cơ sở phân tích nhu cầu và yêu cầu ở trên, đồ án đặt mục tiêu mang lại bốn đóng góp
chính. \emph{Thứ nhất}, một kiến trúc tham chiếu đầu cuối cho giám sát Zero-Trust nội cụm,
kết nối liền mạch từ tầng nhân (eBPF) qua tầng cấu trúc dữ liệu (Count-Min Sketch, đồ thị)
tới tầng dịch vụ và giao diện vận hành --- lấp đúng khoảng trống ``thiếu kiến trúc ráp
hoàn chỉnh'' đã chỉ ra. \emph{Thứ hai}, việc đặt cấu trúc dữ liệu luồng và giải thuật đồ
thị làm trọng tâm phát hiện, kèm phân tích lý thuyết và ví dụ chạy tay (Chương~\ref{ch:foundations}).
\emph{Thứ ba}, một vòng phản hồi cách ly có người duyệt tận dụng \texttt{NetworkPolicy}
của Calico --- biến phát hiện thành hành động ngăn chặn cụ thể mà không cần thay đổi hạ
tầng mạng sẵn có. \emph{Thứ tư}, một bộ số liệu đánh giá định lượng trên cụm Kubernetes
thật (tỉ lệ phát hiện, báo động giả, chi phí eBPF, độ trễ phát hiện$\rightarrow$phản hồi),
giúp kiểm chứng tính thực tiễn của cách tiếp cận.

\section{Phương pháp nghiên cứu}

Đồ án kết hợp nghiên cứu lý thuyết với phương pháp xây dựng -- đánh giá
(\textit{design science}): không dừng ở phân tích mà hiện thực một hệ thống thật rồi đo đạc
để kiểm chứng giả thuyết. Bốn nhóm hoạt động được tổ chức theo đúng mạch các chương sau.

\textbf{Nghiên cứu lý thuyết.} Phân tích các bài báo gốc của từng giải thuật cốt lõi cùng
các đảm bảo toán học của chúng (giới hạn sai số của Count-Min Sketch, tính đúng và độ phức
tạp của Tarjan, BFS, LPM-Trie), làm cơ sở cho lựa chọn thiết kế ở Chương~\ref{ch:foundations}.

\textbf{Phân tích và thiết kế hệ thống.} Từ tập yêu cầu của chương này, dựng các tạo phẩm
phân tích -- thiết kế theo chuẩn kỹ nghệ phần mềm: sơ đồ use-case và đặc tả, sơ đồ phân rã
chức năng, sơ đồ luồng dữ liệu, sơ đồ tuần tự, ERD và thiết kế giao diện
(Chương~\ref{ch:design}).

\textbf{Cài đặt và kiểm thử.} Hiện thực hệ thống với ngôn ngữ phù hợp từng tầng: C cho
chương trình eBPF, C++17 cho thư viện cấu trúc dữ liệu và pipeline, Go cho các dịch vụ SOC;
kiểm thử đơn vị cho các thành phần then chốt (Chương~\ref{ch:impl}).

\textbf{Triển khai và đánh giá.} Triển khai trên một cụm Kubernetes thật theo mô hình
GitOps, chạy các kịch bản tấn công tự động và đo các chỉ số định lượng --- tỉ lệ phát hiện,
tỉ lệ báo động giả, chi phí eBPF và độ trễ phát hiện$\rightarrow$phản hồi
(Chương~\ref{ch:eval}).

\section{Kết luận chương}

Chương này đã thực hiện trọn vẹn giai đoạn chuẩn bị của một quy trình kỹ nghệ phần mềm
trước khi thiết kế: đi từ bối cảnh, khảo sát nhu cầu thực tế, tới phân tích yêu cầu và
đánh giá khả thi. Từ bối cảnh dịch chuyển kiến trúc một khối sang vi dịch vụ trên
Kubernetes, chương chỉ ra một bề mặt tấn công Đông--Tây mới mà an ninh chu vi truyền thống
gần như không quan sát được. Khảo sát tại một doanh nghiệp vận hành Kubernetes nội bộ quy
mô lớn đã lượng hóa khoảng trống này thành các điểm đau cụ thể --- mà trung tâm là sự
\emph{mù lưu lượng Đông--Tây} của lớp giám sát hướng host hiện tại --- rồi chưng cất thành
tập nhu cầu và, qua phân tích, thành các yêu cầu chức năng và phi chức năng có thể truy
vết. Mô hình mối đe dọa theo chuỗi kill-chain làm rõ vì sao việc phát hiện sớm ở giai đoạn
trinh sát và dịch chuyển ngang là điểm can thiệp hiệu quả nhất.

Đối chiếu với ba bài tổng quan có hệ thống, hai báo cáo công nghiệp và các công cụ
cloud-native hiện có, chương xác định một khoảng trống cụ thể: thiếu một kiến trúc tham
chiếu hoàn chỉnh ráp nối từ tầng nhân tới tầng dịch vụ và giao diện vận hành, lấy cấu trúc
dữ liệu luồng và giải thuật đồ thị làm trọng tâm phát hiện. Trên cơ sở đó, chương đã phát
biểu bốn đóng góp chính, mục tiêu, đối tượng và phạm vi của đề tài. Toàn bộ các yêu cầu và
căn cứ này là đầu vào trực tiếp cho việc phân tích, cụ thể hóa thành thiết kế ở
Chương~\ref{ch:design}, sau khi đã trang bị cơ sở lý thuyết và công nghệ lõi ở
Chương~\ref{ch:foundations}.
